\documentclass[border=7bp]{standalone}

\usepackage{fontspec}
\usepackage{xeCJK}
\setmainfont{Calibri}
\setCJKmainfont{SimSun}[AutoFakeBold=1.5]
\usepackage{amsmath}
\usepackage{tikz}
\usetikzlibrary{arrows.meta,shapes.geometric}

\definecolor{Ink}{HTML}{25313D}
\definecolor{Blue}{HTML}{315E80}
\definecolor{PaleBlue}{HTML}{E7EFF5}
\definecolor{Muted}{HTML}{626B74}
\definecolor{PaleGray}{HTML}{F0F1F3}

\tikzset{
  >=Stealth,
  every node/.style={font=\fontsize{9bp}{12bp}\selectfont,text=Ink},
  edge/.style={->,draw=Ink,line width=.6pt},
  box/.style={rectangle,draw=Ink,line width=.6pt,align=center,
              text width=6.0cm,minimum height=0.95cm,inner sep=3pt},
  decision/.style={diamond,aspect=3.0,draw=Ink,line width=.6pt,align=center,
                   text width=3.6cm,inner sep=1pt},
  side/.style={rectangle,draw=Muted,dashed,line width=.65pt,align=center,
               fill=PaleGray,text width=3.4cm,inner sep=4pt,
               font=\fontsize{8.5bp}{11bp}\selectfont},
  small/.style={font=\fontsize{8.5bp}{11bp}\selectfont},
}

% Shared preamble for the report's editable vector figures.
\usepackage{fontspec}
\usepackage{xeCJK}
\setmainfont{Calibri}
\setCJKmainfont{SimSun}[AutoFakeBold=1.5]
\usepackage{amsmath}
\usepackage{tikz}
\usetikzlibrary{arrows.meta,positioning,calc,shapes.geometric,fit}
\definecolor{Ink}{HTML}{25313D}
\definecolor{Blue}{HTML}{315E80}
\definecolor{PaleBlue}{HTML}{E7EFF5}
\definecolor{Muted}{HTML}{626B74}
\definecolor{PaleGray}{HTML}{F0F1F3}
\tikzset{
  >=Stealth,
  every node/.style={font=\fontsize{10bp}{13bp}\selectfont,text=Ink},
  flow/.style={->,draw=Ink,line width=.65pt},
  frame/.style={rectangle,draw=Ink,line width=.65pt,align=center,
    minimum height=.85cm,text width=5cm,inner sep=5pt},
  emphasis/.style={frame,draw=Blue,fill=PaleBlue},
  supporting/.style={frame,draw=Muted,fill=PaleGray},
  annotation/.style={font=\fontsize{9bp}{12bp}\selectfont,text=Muted,align=center},
}

\begin{document}
\begin{tikzpicture}
  \node[box] (start)  at (0,0)     {起点：抽象目标状态，距离置 0 入队};
  \node[box] (pop)    at (0,-2.05) {从\textbf{队首}取出 $c$，还原位置};
  \node[box] (cost)   at (0,-4.10) {考察空位四邻：换入本组成员\\则 $w=1$，否则 $w=0$};
  \node[decision] (better) at (0,-7.20) {候选 $c+w$\\小于现值？};
  \node[box] (write)  at (0,-10.30) {写入该值，决定从队头或队尾进};
  \node[box] (queue)  at (0,-12.35) {$w=0$：\texttt{addFirst}（队头）\\$w=1$：\texttt{addLast}（队尾）};

  \node[box,text width=2.4cm] (skip) at (5.40,-7.20) {不做改动};

  % 哨兵注解移到右上角，避开左侧回边竖线
  \node[side] (sentinel) at (5.40,-3.10)
    {\textbf{哨兵约定}\\\texttt{0xFF} 即 255 未访问\\合法距离 $0$--$254$\\候选 $\ge255$ 即报错};

  \node[side] (cache) at (5.40,-1.00)
    {\textbf{磁盘缓存}\\\texttt{generated/pdb/}\\魔数、版本、规模、组员、\\目标指纹、CRC32 全符\\则载入，否则建表写回};

  \draw[edge] (start.south)  -- (pop.north);
  \draw[edge] (pop.south)    -- (cost.north);
  \draw[edge] (cost.south)   -- (better.north);
  \draw[edge] (better.south) -- node[right,small]{是} (write.north);
  \draw[edge] (write.south)  -- (queue.north);
  \draw[edge] (better.east)  -- node[above,small]{否} (skip.west);

  \draw[edge] (queue.west) -- (-4.30,-12.35)
    -- node[left,small,pos=.55]{队列非空} (-4.30,-2.05)
    -- (pop.west);

  % 显式声明边界，确保左右两侧注解框不被裁掉
  \path[use as bounding box] (-6.60,0.75) rectangle (7.45,-13.20);
\end{tikzpicture}
\end{document}
